\section{Introduction}
\label{sec:intro}
The constitutive response of a heterogeneous solid is an outcome of the interaction between its constituents, their spatial arrangement, and the deformation imposed on them. At finite strains, this interaction can produce anisotropy and couplings between loading modes that are difficult to represent with a small, predetermined set of macroscopic material parameters. Computational homogenization addresses this difficulty by replacing the explicit macroscopic constitutive equation with a microscopic boundary-value problem. The macroscopic deformation drives a representative volume element (RVE), and the equilibrated microscopic fields determine the effective stress and tangent. The formulation of Miehe et al.~\cite{Miehe1999}, developed for finite crystal plasticity, illustrates this micro-to-macro transition and distinguishes uniform deformation, uniform boundary deformation, and periodic fluctuation constraints. In the present work, the microscopic material is hyperelastic and the coupling is periodic.

The resulting two-level finite element method, commonly denoted FE$^2$, is attractive because it retains a direct connection between the effective response and the underlying microstructure. Its computational cost, however, increases rapidly: each macroscopic integration point carries a microscopic problem, and the latter may have to be solved repeatedly during a macroscopic Newton iteration. This cost is particularly consequential in a many-query setting, where a fixed material is used in several structural analyses, optimization iterations, or parameter studies. The question is therefore not simply how to evaluate a fitted stress--strain curve quickly, but how much of the microscopic model must remain online to achieve a prescribed accuracy and retain the mechanical properties required by the macroscopic calculation.

Projection-based model order reduction (PMOR) and learned constitutive modeling offer two different answers. PMOR reduces the dimension of the microscopic state while retaining equations inherited from microscopic equilibrium. Constitutive learning instead approximates the effective input--output relation and can dispense with microscopic equilibrium altogether. Both use information generated offline, and both introduce approximation error, but they approximate different mathematical objects. A neural network that predicts a microscopic closure is not the same model as a neural network that predicts the effective energy. Keeping this distinction explicit is essential when comparing their accuracy, mechanical guarantees, and computational cost.

An early application of reduced microscopic kinematics to finite-strain homogenization is the reduced model multiscale method of Yvonnet and He~\cite{yvonnet2007}. Subsequent developments address both approximation and computational efficiency. Hern\'andez et al.~\cite{Hernandez2014} construct reduced microscopic spaces and examine the stability of the additional approximation used to evaluate nonlinear terms. Fritzen and Hodapp~\cite{Fritzen2016} combine a potential-based reduced homogenization model with GPU acceleration for three-dimensional two-scale calculations. These contributions establish that reduced micromechanics is a substantive alternative to full nested analysis, rather than merely a convenient baseline for a learned constitutive law. The domain-decomposition method of Hern\'andez~\cite{Hernandez2020} also demonstrates a different use of reduced periodic substructures, including interactions between adjacent cells. That approach should be distinguished from the scale-separated, first-order FE$^2$ setting adopted here.

The conventional projection-based reduced-order model (PROM) approximates the high-dimensional solution in an affine subspace. Proper orthogonal decomposition (POD) provides an optimal approximation to the training snapshots in a specified discrete norm, but it does not ensure that a very small affine space represents a nonlinear solution manifold accurately. Several strategies address this limitation. Local reduced-order bases partition the relevant solution space~\cite{Amsallem2012}; nonlinear manifold approaches include convolutional autoencoders~\cite{Lee2020} and quadratic approximation manifolds~\cite{Barnett2022}. The PROM--ANN formulation of Barnett et al.~\cite{barnett2023} instead augments a primary reduced basis with a learned closure in a secondary reduced space. Ares De Parga et al.~\cite{aresdeparga2026nonlinear} develop this latent-space closure viewpoint further and compare neural, Gaussian-process, and radial-basis-function approximations. Their notation and distinction between primary generalized coordinates, secondary generalized coordinates, and the closure map are adopted throughout the present paper.

This distinction has a practical consequence. A nonlinear PROM can represent displacement patterns in a comparatively large POD space while solving online for only a few primary generalized coordinates. It does not need a neural network whose output dimension equals the high-dimensional state size. However, fewer coordinates do not by themselves eliminate the cost of evaluating the microscopic residual. Hyperreduction is needed to reduce that second cost. Optimized cubature for subspace deformations~\cite{An2008}, energy-conserving sampling and weighting (ECSW)~\cite{Farhat2014,farhat2015}, and the empirical cubature method (ECM)~\cite{hernandez2017} approximate the relevant reduced integrals using selected contributions and weights. Extensions address Petrov--Galerkin projection and local bases~\cite{Grimberg2021,AresDeParga2023}. These methods make the distinction between a low-dimensional system and a computationally inexpensive system concrete.

The cubature construction itself can also adapt to the approximation space. Continuous empirical cubature adjusts point locations as well as weights~\cite{Hernandez2024}, while subspace-adaptive cubature uses common locations with weights tailored to different reduced subspaces~\cite{Bravo2024}. Manifold-adaptive empirical cubature extends the adaptive-weight viewpoint to nonlinear approximation manifolds~\cite{hernandez2026}. The present study uses fixed element-selected ECM for its linear HPROM and manifold-adaptive weights for its HPROM--ANN. Each is deployed on an actual reduced finite element mesh. These choices are established reduction ingredients; the contribution here is their explicit integration and assessment as competing microscopic constitutive evaluators, with their costs and assumptions made comparable.

Learned macroscopic constitutive laws remove another layer of online computation. Their diversity reflects the variety of quantities that can be learned: homogenized stresses from microstructural information~\cite{Aldakheel2023}, reduced representations of fine-scale histories~\cite{Liu2022}, or energetically structured effective responses~\cite{asad2022,kalina2023}. The staggered multiscale framework of As'ad and Farhat~\cite{AsadFarhat2026} further addresses the offline cost of learning across several interacting scales in path-dependent woven materials. These developments motivate constitutive surrogates for expensive multiscale calculations, but they do not make the mechanics of every surrogate interchangeable. The present investigation isolates a history-independent material and compares direct stress regression, an unrestricted scalar-energy network, and two polyconvex energy constructions.

The first distinction is the existence of an energy potential. As'ad et al.~\cite{asad2022} show why fitting the components of stress independently does not enforce the integrability required of a hyperelastic law. Predicting a scalar energy and differentiating it enforces this relation and the corresponding tangent symmetry. Their use of an input-convex neural network (ICNN) in Green--Lagrange strain additionally constrains strain-space curvature. However, convexity in this strain measure and polyconvexity in the deformation gradient are different requirements. Gao et al.~\cite{GaoNeffRoventaThiel2017} clarify the role of convexity in the right Cauchy--Green tensor and of the second Piola stress in associated minimization arguments. In particular, strain-space convexity cannot be substituted for a general finite-strain rank-one-convexity condition.

Polyconvexity, introduced by Ball~\cite{ball1977}, expresses the energy as a convex function of deformation minors. Together with appropriate growth, coercivity, and admissibility hypotheses, it supports existence results in nonlinear elasticity~\cite{Ciarlet1988}; it is neither a uniqueness theorem nor a guarantee that a particular Newton algorithm converges. In anisotropic constitutive modeling, invariant formulations using structural tensors provide a systematic way to encode directional information. Schr\"oder and Neff~\cite{SchroderNeff2003} develop polyconvex invariant building blocks for transverse isotropy and discuss an orthotropic extension. Gasser et al.~\cite{Gasser2006} illustrate how an additional structural description of fiber dispersion changes the predicted deformation mechanisms. Their example is relevant to the broader modeling issue: an admissible functional form can still be too restrictive to represent the material mechanisms of interest.

Neural architectures extend these ideas by learning the combination of admissible constitutive ingredients. Klein et al.~\cite{klein2022} develop invariant-based and deformation-gradient-based polyconvex anisotropic networks, with distinct treatments of objectivity. Linden et al.~\cite{linden2023} formulate compatible growth and reference-normalization terms, while treating nonnegative energy as a separate issue. Constitutive artificial neural networks (CANNs)~\cite{Linka2023} and polyconvex neural ordinary differential equations~\cite{Tac2022} offer other parameterizations of admissible energies. The benchmarking study of Tac et al.~\cite{Tac2024} is particularly pertinent: close agreement in stress does not necessarily imply similar second derivatives. Generalized-structure-tensor networks~\cite{Kalina2024} and objective triclinic polyconvex models~\cite{klein2026} also address anisotropy beyond a prescribed, simple invariant set. Accordingly, learning directions without prescribing a material symmetry class is not, by itself, a novelty claim of the present work.

Related terminology can conceal important differences. The positive-matrix architecture of Xu et al.~\cite{Xu2021} concerns an incremental constitutive update; a symmetric positive semidefinite matrix factor is not by itself a hyperelastic potential. The directional learning-agent construction of Ghaderi et al.~\cite{Ghaderi2020} incorporates microsphere structure and history variables to describe polymer inelasticity. Constitutive Kolmogorov--Arnold networks (CKANs)~\cite{Abdolazizi2025} combine energy modeling, partial monotonicity, and symbolic interpretability, whereas input-convex Kolmogorov--Arnold networks (ICKANs)~\cite{thakolkaran2025} additionally constrain the spline-based core to construct polyconvex energies. An assessment of these approaches must specify the modeled state, the variables in which convexity is imposed, and the constraints that the particular architecture actually enforces.

Within this established landscape, an important practical difficulty remains at the interface between admissibility and representation. Klein et al.~\cite{klein2022} already provide a revealing example: their invariant model struggles with shear in a cubic lattice, while their deformation-gradient model fits it more accurately; the invariant approach performs well on their transversely isotropic example. This contrast makes the feature representation, rather than polyconvexity alone, a central modeling choice. A monotone convex network can only fit responses compatible with the ordering induced by its features. If a prescribed feature vector increases componentwise along a deformation path, the network core cannot decrease along that path. Additional normalization terms must be included before drawing a conclusion about the complete energy, but increasing the core's width does not remove an incompatible feature ordering. Conversely, discarding all constraints may improve approximation while forfeiting a useful mechanical guarantee. The relevant design question is therefore how to enlarge the admissible representation without abandoning the properties that motivated it.

This work addresses that question with paired directional features whose powers, orientations, and determinant exponents are trainable within explicit convexity bounds. Pairing gives each feature an isotropic derivative at the reference state even though its finite-deformation response can be anisotropic. Consequently, an affine determinant correction can enforce a stress-free reference independently of the positive volumetric barrier. The same feature construction is used with a smooth ICNN and a convex spline ICKAN, so that their comparison concerns two admissible core parameterizations within a common constitutive design. A separate sufficient lower-bound argument examines nonnegative energy for the selected trained parameters.

The numerical assessment is organized around two complementary questions. The periodic-cell study asks whether the constrained representation can fit the homogenized anisotropic response and which mechanical distinctions are detectable beyond pointwise stress error. A closed strain cycle tests integrability; directional second variations test local rank-one curvature. The structural coupon study then asks how the constitutive approximations and the reduced microscopic models propagate into finite element errors and online wall time. This division follows the principle that an example should resolve a specific question: the coupon is a structural deployment of the same material, not a second independent material validation. The need to accompany global errors with interpretable field comparisons also follows the lesson of the reduced-model assessment in Chmiel et al.~\cite{Chmiel2024}, although the physical problem and error measures here are different.

The paper thus contributes an admissible anisotropic feature construction, its neural and spline realizations, and a common assessment against optimized reduced micromechanics and nested full-order homogenization. Figure~\ref{fig:hierarchy} summarizes the two approximation routes and their common structural interface. The evidence establishes useful accuracy and efficiency for one microstructure and a specified deformation domain. It does not require an assertion of universal approximation, universal extrapolation, or automatic superiority over unrestricted models. Section~\ref{sec:fom} defines the reference homogenization problem. Section~\ref{sec:reduced} develops the reduced microscopic models, and Section~\ref{sec:learned} presents the learned constitutive laws. Sections~\ref{sec:cell} and \ref{sec:coupon} address constitutive behavior and structural deployment, respectively. The final discussion relates the measured trade-offs to the mathematical scope of each model.

\begin{figure}[htbp]
 \centering\includegraphics[width=.98\linewidth]{model_hierarchy.pdf}
 \caption{Two routes from offline microscopic information to an online macroscopic constitutive evaluation. A learned microscopic closure and a learned effective energy are different approximation objects. Only the iterative reduced microscopic route retains an online microscopic equilibrium solve. The diagram describes the implemented workflow; it is not a statement of equal training cost or equal mechanical guarantees.}
 \label{fig:hierarchy}
\end{figure}
